% !TEX root = ../main.tex
% Figure: FDTD simulation domain setup for PCSEL
% Chapter: ch14 - FDTD and FEM Methods for PCSEL Simulation
% Description: Schematic of a typical 3D FDTD simulation setup

\begin{figure}[htbp]
  \centering
  \begin{tikzpicture}[scale=1.2]
    % Main simulation volume (air box)
    \draw[thick, blue!60!white] (-3,-2) rectangle (3,3);
    \fill[blue!5] (-3,-2) rectangle (3,3);
    \node[above right, blue] at (-3,3) {FDTD 计算区域};
    
    % PML absorbing boundaries
    \draw[line width=3mm, gray!50] (-3.15,-2.15) rectangle (3.15,3.15);
    \node[right, gray] at (3.2,2.5) {PML 吸收层};
    \node[font=\small, gray] at (3.8,2.5) {\scriptsize 12--20 层};
    
    % PCSEL device structure
    \fill[green!20] (-2,-1.5) rectangle (2,-0.5);
    \draw[thick] (-2,-1.5) -- (2,-1.5);
    \draw[thick] (-2,-0.5) -- (2,-0.5);
    \node[right] at (2.1,-1) {PCSEL 有源层};
    
    % Photonic crystal holes
    \foreach \x in {-1.5,-0.75,...,1.5} {
      \foreach \y in {-1.25,-0.75} {
        \fill[white] (\x,\y) circle (0.15);
        \draw[thick] (\x,\y) circle (0.15);
      }
    }
    \node[above] at (0,-0.5) {PhC 孔洞阵列};
    
    % Dipole source inside active region
    \fill[red] (0,-1) circle (3pt);
    \draw[->, very thick, orange] (-0.3,-1) -- (0.3,-1);
    \node[below, orange] at (0,-1.3) {偶极子源};
    \node[font=\small, orange] at (0,-1.55) {\scriptsize $\mathbf{J}(t)=\mathbf{J}_0 e^{-i\omega t}$};
    
    % Field monitors
    \draw[dashed, purple, thick] (-1.5,-2) -- (-1.5,3);
    \node[purple, rotate=90] at (-1.6,0) {时间监视器};
    
    \draw[dashed, magenta, thick] (-3,0) -- (3,0);
    \node[magenta] at (2.5,0.3) {频谱监视器};
    
    % Near-field to far-field transformation surface
    \draw[dashed, cyan, thick] (-2.5,-1.8) rectangle (2.5,2.5);
    \node[cyan] at (0,2.7) {近远场变换面};
    
    % Mesh grid overlay
    \foreach \x in {-2.5,-2,...,2.5} {
      \draw[gray!20, very thin] (\x,-2) -- (\x,3);
    }
    \foreach \y in {-1.5,-1,...,2.5} {
      \draw[gray!20, very thin] (-3,\y) -- (3,\y);
    }
    \node[below right, gray] at (-3,-2) {Yee 网格};
    
    % Time evolution inset
    \begin{scope}[shift={(4.5,1)}, scale=0.8]
      \draw[->] (0,0) -- (3.5,0) node[right] {$t$};
      \draw[->] (0,0) -- (0,2.5) node[above] {$E_z$};
      
      % Pulse excitation
      \draw[thick, red, domain=0:3, samples=100] 
        plot(\x,{1.5*exp(-(\x-1)*(\x-1)/0.1)*sin(deg(20*\x))});
      
      \node[above, font=\small] at (1.5,2.7) {脉冲激励};
      \node[below, align=center, font=\footnotesize] at (1.5,-0.3) {宽带覆盖\\$\Delta\lambda \sim 100\,\mathrm{nm}$};
    \end{scope}
    
    \begin{scope}[shift={(4.5,-1)}, scale=0.8]
      \draw[->] (0,0) -- (3.5,0) node[right] {$t$};
      \draw[->] (0,0) -- (0,2.5) node[above] {$|E|^2$};
      
      % Ring-down decay
      \draw[thick, blue, domain=0:3.2, samples=100] 
        plot(\x,{2*exp(-\x/0.8)*(1+0.3*sin(deg(15*\x)))});
      
      \node[above, font=\small] at (1.5,2.7) {指数衰减};
      \node[below, align=center, font=\footnotesize] at (1.5,-0.3) {拟合得$Q$值\\$\tau = Q/\omega_0$};
    \end{scope}
    
    % Annotations
    \draw[<->, thick] (-3.5,-2) -- node[left] {$L_z$} (-3.5,3);
    \draw[<->, thick] (-3,-2.5) -- node[below] {$L_x$} (3,-2.5);
    
    % Legend
    \node[align=left, font=\small, anchor=north west] at (-3,-2.3) {
      \textbf{典型设置}:\\
      $\bullet$ 网格尺寸:$\Delta x < \lambda/(10n)$\\
      $\bullet$ 时间步长:CFL 条件$\Delta t < \Delta x/(c\sqrt{3})$\\
      $\bullet$ 总时长：数万个时间步直至充分衰减
    };
  \end{tikzpicture}
  \caption{三维 FDTD 仿真的典型配置示意图。蓝色方框为主计算区域（通常为空气背景），外围灰色厚边框代表完美匹配层 (PML)——一种人工吸收边界，用于模拟开放空间并抑制虚假反射。绿色矩形为 PCSEL 的外延层截面，其上刻蚀有周期性排列的空气孔（白色圆圈）。红色圆点表示放置在有源区中心的偶极子脉冲源$\mathbf{J}(t)$，其宽带频谱可以同时激发多个谐振模式。紫色竖线和品红色横线分别代表时域和频域监视器，用于记录指定位置的场随时间的演化和稳态频谱响应。青色虚线框为近远场变换 (Jordan-Balakrishnan) 的闭合表面积分面，通过记录该面上的切向场分量可以重建任意远处的辐射方向图。灰色的 Yee 网格展示了 FDTD 的空间离散方案——电场和磁场分量交错布置在每个元胞的棱边和面上，保证了 Maxwell 方程的守恒性质。右上角插图显示了典型的脉冲激励波形（高斯包络正弦波）和随后的指数衰减尾迹，从中可以提取品质因子$Q=\omega_0\tau/2$。左下角总结了经验法则：为保证数值色散可忽略，网格尺寸应小于介质中波长的 1/10；时间步长受 CFL 稳定性条件限制。}
  \label{fig:ch14_fdtd_setup}
\end{figure}
